\section{边疆与南方：王室的地图走到哪里，能力就走到哪里}

一支王师若从关中出发，越过成周所联结的东方网络，再向汉水方向推进，带着的不只是戈车和旗帜。粮食要沿途接济，车马要有可走的道路，随行的贵族与军士要在陌生的水土中维持队伍；更要紧的是，沿线的政治共同体未必把周王的命令当作自己的命令。西周王室的疆界因而不是一根已经划定的线，而是补给、交通、盟友与威慑能够反复抵达的距离。

\phantomsection\label{fig:vol01:western-zhou:zhao-southern-routes}
\begin{center}
\begin{tikzpicture}[x=.88cm,y=.88cm,font=\small]
  \fill[paper] (0,0) rectangle (11.2,6.0);

  \draw[source!65!timeline,line width=1pt]
    (6.7,.35) .. controls (7.4,1.4) and (7.0,2.6) .. (8.2,3.55)
    .. controls (9.1,4.25) and (9.0,5.0) .. (10.5,5.65);
  \node[text=source!80!ink,right] at (8.35,2.2) {汉水};
  \draw[source!42,dashed] (3.1,5.15) .. controls (4.2,4.65) and (5.1,4.7) .. (5.9,4.2);
  \node[text=source!80!ink,above] at (4.5,4.85) {秦岭南麓方向};

  \fill[dynasty!18] (.7,3.75) -- (2.7,3.55) -- (3.0,4.85) -- (1.0,5.15) -- cycle;
  \node[align=center] at (1.85,4.35) {宗周\\关中};
  \fill[timeline!20] (3.95,2.7) -- (5.85,2.55) -- (6.15,3.85) -- (4.25,4.1) -- cycle;
  \node[align=center] at (5.05,3.3) {南阳—汉水\\北部走廊};
  \fill[source!14] (7.35,.65) -- (10.35,.75) -- (10.55,2.65) -- (8.1,2.95) -- cycle;
  \node[align=center] at (9.1,1.75) {汉水以南\\接触区域};

  \draw[-{Stealth[length=2mm]},ultra thick,color=dynasty]
    (2.8,4.2) -- (3.78,3.65);
  \draw[-{Stealth[length=2mm]},ultra thick,color=timeline]
    (6.2,3.05) .. controls (6.9,2.55) and (7.35,2.45) .. (8.0,2.2);
  \node[text=dynasty,above,align=center] at (3.15,4.05) {王室调动\\与补给};
  \node[text=timeline,above,align=center] at (7.05,2.85) {军旅、物资与\\信息链延伸};

  \draw[-{Stealth[length=2mm]},thick,dashed,color=source]
    (8.3,2.7) .. controls (9.65,3.2) and (10.1,3.9) .. (9.55,4.6);
  \node[draw=dynasty!75,fill=paper,rounded corners=2pt,align=center,inner sep=4pt]
    at (9.35,5.0) {汉水失利的\\传统叙事所系\\地点未定};
  \node[text=source,align=center,font=\footnotesize] at (4.3,.85)
    {接触可由器物与考古旁证；\\不应还原为一条连续、确定的远征路线};
\end{tikzpicture}

\smallskip
\small\textit{图\quad 昭王南向行动与汉水风险示意：用宗周至汉水以南的交通、补给与信息距离，说明多次南向接触和后来“丧师汉水”叙事所面对的空间约束。参考《史记·周本纪》、今本《竹书纪年》、《左传·昭公十三年》以及南方青铜器、汉水流域考古；行动次数、目标、行军路径和失利地点均有争议。所有边界与路线均为约略示意。}
\end{center}


\begin{event}{西周中期（具体王年未定）}{器物到达南方，不等于王命已经抵达}{器物与接触可靠；路线、行动对象和政治关系有争议}{汉水、江汉及东南方向的多条水陆通道；王室、贵族与地方政治体}\eventanchor{WZ-MID-SOUTH-ROUTE-INSCRIPTIONS}{西周·南方通道与器物接触}
一件周式铜器在汉水或江汉方向出现，首先说明这片区域并非王室视野以外的空白，却不替我们说明它为何到达。赠赐、交换、婚姻、迁徙、掠夺和受命行动都可能留下器物；器物分布所见的是持续接触的条件，而非一条由宗周单向延伸、每一站都受控的道路。对王室来说，铜料、河道和地方联盟可能共同构成南向经营的诱因；对沿线共同体来说，与周发生联系也并不必然意味着服从。

若把器物读回其铭文，又能看见远方行动需要许多中介。小臣单觯记近臣参与军旅或俘获，麦方尊和静方鼎各保存贵族处理地域军政事务、以功受赏的线索。它们的对象、地望和彼此关联都有异说，不能拼接为一支部队从关中抵达江汉的完整路线；但其共同之处颇为清楚：王命必须经由近臣、车马、受命贵族和地方关系才会成为行动。战争和接触因而是多点发生的过程，而不是中心在地图上画出一道箭头。

淮夷方向的金文亦使这种判断延伸到东南。王室与贵族的征伐、防御和赏赐可以被可靠地看见，连续战争的年次和具体地望却不能确定。周初的东方安置没有把这些区域变成无摩擦的腹地；南方与东南的道路越长，需要传递的粮食、命令和信息便越多。后来\eventref{WZ-0977-ZHAONAN}{西周·昭王南征}的失败，正应放在这一片已经存在而又始终不稳的接触网络中理解。
\end{event}

\begin{event}{昭王在位后段（约前10世纪，具体年次未定）}{周昭王南征失利}{王年推定}{汉水流域及其以北、以东的南向通道}\eventanchor{WZ-0977-ZHAONAN}{西周·昭王南征}
昭王把王室的行动推向南方，并非一次偶发的冒险。早期西周经由\eventref{WZ-CHENGKANG-CONSOLIDATION}{西周·成康巩固}所形成的册命、封国和征调关系，给王室提供了远征的条件，也给它提出了继续展示支配力的要求。南方的金属资源和交换路线、沿江汉通道的贡纳与服从、以及在新近接触的共同体面前维持王命的声望，都可能是这类行动的使能条件；现存材料却没有一篇出征诰命，容许我们把其中任何一项写成昭王唯一而明言的目标。所谓“南方战略”，应当理解为多次经营、征伐和关系重组的方向，不是一张已经完成的征服地图。

传世叙事把这条路的终点放在汉水。\sourcebook{史记·周本纪}只说昭王“南巡狩不返，卒于江上”；今本\sourcebook{竹书纪年}则有“丧六师于汉”的浓缩记忆；\sourcebook{左传·昭公十三年}中楚人以“昭王南征而不复”反诘周使。三种说法共享一个坚硬的核心：周王曾在南向行动中未能归来，王室遭受了严重损失。它们并不提供一份同时代战报。尤其“六师”既可能保存王师覆败的规模感，也可能是后出叙事以王室军制名目作出的概括，不能据此复原六支部队各自的编制、渡河地点和死伤数字。

灾难究竟发生在昭王在位的哪一年，同样不能用整齐的年表遮蔽。昭王的绝对年代本来依赖后世编年与现代推算；《竹书纪年》文本的传承和异文又使其中年次难作不移的坐标。汉水的水难、战斗失利，还是二者在撤退中相继发生，传世材料也没有足够细节裁决。可以确认的是，渡河这一具体的地理障碍进入了失败记忆；不能确认的，是王师在何处列阵、由谁击溃，或昭王的遗体如何处置。把不确定的河岸写成一处确知战场，反而会削弱这场失败所留下的真实重量。

后世的楚国很容易使这次远征看上去像周、楚两国之间的一场定型战争。那是倒读。昭王时代的江汉并非一个由单一“楚国”统辖的南方，楚的早期政治共同体正在形成，汉水与江汉之间还分布着许多族群、聚落和具有自身权力的首领网络。楚人后来保存并使用昭王不返的记忆，说明这场失败进入了周楚关系的政治语言；却不足以证明昭王所面对的每一支力量都是后来的楚，或一次战役已经决定了楚的崛起。南方诸政治体的成长与周王室的推进、交换和冲突彼此缠绕，不能压缩成一方突然取代另一方的故事。

考古材料使这种谨慎更有必要。江汉、汉水流域的西周遗址、墓葬和青铜器，能够看见周式礼器、技术与地方传统的交错，也能提示这片区域绝非王室视野以外的空白；它们不能指出一处“丧六师”的战场，更没有一件已知铭文直书昭王在汉水覆亡。器物的流通可以来自赠赐、贸易、婚姻、迁徙或武力，墓地的文化面貌也不能直接等同于一支军队的行军线。考古学在这里提供的是南方社会复杂而有联系的背景，而不是为后出的叙事补齐战斗场景。

昭王未归，首先是一次人事与指挥的断裂：王位继承必须在王室失去君主、远征队伍受挫的处境中处理，幸存的贵族和承担征调的政治网络也要面对损失。若王师确有大规模伤亡，车马、甲兵、役夫和供给的缺口便不会随着一次册命自动消失；王室还须重新分配赏赐、劳役和军事义务，才能维持对东方与边疆的关系。没有材料能够证明此后立即实行了某项统一军制改革，但这场失败把王权远征的成本暴露给了朝廷：命令能够发出，不等于补给和归途也能由中心掌握。

因此，\eventref{WZ-0977-ZHAONAN}{西周·昭王南征}不应被当作西周覆亡的单一开端。它留下的是一项更长久的约束：王室在动员远方时，必须把交通风险、地方联盟和军费一并计入。到\eventref{WZ-XUAN-REVIVAL}{西周·宣王中兴}所面对的料民与边防压力时，问题已转到西北，材料也不同；两者之间不能画一条直接的因果直线。不过，南征的惨重代价已提醒人们，维持一张跨区域的王权网络，收益和损耗总会一同回流到宗周。
\end{event}

\begin{event}{厉王时期说（约前9世纪；具体王年未定）}{鄂侯驭方之争：胜方铭文留下了什么}{金文所记冲突可靠；人物身份、原因、范围与“叛”字解释有争议}{南方相关区域；王室、鄂侯驭方及受命贵族}\eventanchor{WZ-LI-EHOU-YU}{西周·鄂侯驭方冲突}
禹鼎和鄂侯驭方器群把一次冲突的痕迹留在不同的器面上。前者所见军政调动和赏赐，使王室集中军力、酬报服务者的能力仍然可见；后者则保存与鄂侯驭方有关的另一组线索。器物存在并不等于事件的每一个环节都已明白：参与者的具体地位、战区的大小、资源冲突还是政治服从何者更为直接，以及某些关键字是否应释作“叛”，都仍须谨慎。

正因如此，这场冲突不宜被写成南方已经形成统一反周阵线的证据。较稳妥的说法是，王室、区域贵族和地方政治体各自拥有可调动的资源，关系破裂时需要以军事和赏赐重新排列。禹鼎所见的胜利语言首先服务于功劳的确认，而非替失败者陈述理由；与鄂侯驭方器并读，恰好提醒读者不要把单一胜方叙述扩展成完整的政治地理。它所显示的，是晚西周南方控制仍须反复经过贵族军事中介，而不是一项已经完成的行政接管。
\end{event}

\begin{event}{宣王时期说（具体王年、战役与封地性质未定）}{北方在用兵，江汉与淮夷方向也没有静下来}{诗歌与金文保存军政记忆；作者、对应战役和人物行迹有争议}{江汉、淮河及申、谢相关地区；王室、受命贵族与地方政治体}\eventanchor{WZ-XUAN-HUAIYI}{西周·宣王时期的南方军政行动}
\sourcebook{诗经·大雅·江汉}和\sourcebook{常武}用颂扬性的语调保留南方、东南军事行动的记忆。它们能说明王室仍愿把对边缘政治体的处置写成恢复秩序的功业，却不能让我们据篇章逐句复原一次战役的日期、将领和行军方向。东部金文与淮河流域考古所见的反复接触，可以与诗歌构成趋势上的互证；“淮夷”所指的群体和每次行动的对象仍不清楚，不能将其当成固定不变的敌国名。

\sourcebook{崧高}所叙申伯受命赴谢、营邑赐土，则把军事之外的安排带到眼前。诗篇使营邑、宗族和王命相连，谢地的性质和具体经营过程却有争议；它较适合显示王室要在南方交通与联盟地带设置可依托的贵族支点，不能证明一座城或一块封地从此完全由中心直接治理。与西北的征伐、料民并列来看，宣王时期的困难不在于有没有一场可歌咏的胜利，而在于不同方向同时需要人马、供给、受命者和地方合作。

这也为《江汉》《常武》的胜利声调增加了一层重量。南北并行的动员可以暂时重申王命，却会提高财政与协调的要求；这一点是从材料的并置中作出的结构性解释，不能换算成确切兵额或税额。南方行动没有直接造成\eventref{WZ-0771-HAOJING}{西周·镐京陷落}，但它与西北边防一样，说明西周晚期的王权必须不断把远方的压力带回宗周处理。
\end{event}

\begin{sourcenote}[一条河与三层材料]
\sourcebook{史记·周本纪}、\sourcebook{竹书纪年}与\sourcebook{左传·昭公十三年}成文、传抄和叙事目的各不相同：前者是汉代本纪，后两者保存的也是后出文本层中的历史记忆。它们可以互证昭王南行不返这一叙事核心，不能互相补足为一份精确的作战日志。江汉地区西周考古与早期楚研究则应同这些文献并读：它们校正“南方空白”与“单一楚国”的想象，却尚未提供昭王覆师的直接物证。
\end{sourcenote}

\begin{sourcenote}[铭文、器物分布与颂诗不能互相替代]
小臣单觯、麦方尊、静方鼎、禹鼎和鄂侯驭方器群所提供的是器主关切的受命、功劳、俘获或冲突线索；铭文可靠，不表示其中的对象、地望和动机都可无争议地还原。器物分布与墓葬材料能够校正“南方空白”的想象，却不能直接证明某一支王师经过何处。《江汉》《常武》《崧高》保存的则是颂扬、营邑和王命的政治语言。三类材料并读，可以辨认南方经营的多重成本；将它们顺次拼成完整战报或精确疆界，反会超过它们各自所能承担的证据。
\end{sourcenote}
